\section{Conclusion}
\label{sec:conclusion}

We presented a systematic study of character tracking capacity in frontier language models, varying the number of characters from 2 to 64 in controlled synthetic narratives.
Our key findings are: (1)~both \gptfull and \gptmini achieve perfect accuracy up to 16~characters and maintain above 93\% accuracy at 64~characters; (2)~when errors occur, they are overwhelmingly confusions between characters (93.9\% for \gptmini), not forgetting; (3)~degradation is gradual, with no sharp cliff; and (4)~the smaller \gptmini outperforms the larger \gptfull, partly due to a specific failure mode in \gptfull on compound possession-change sentences.

These results establish a new empirical baseline for entity tracking in frontier models and provide error-pattern evidence consistent with the associative tracking mechanisms identified by recent interpretability work~\citep{li2025state}.
The confusion-dominated error pattern suggests that character attributes are stored in overlapping distributed representations, where increased entity count leads to cross-character interference rather than information loss.

\para{Future work.}
Several directions follow naturally from this work.
Testing additional model families (Claude, Gemini, open-weight models) would establish the generality of our findings.
Extending to naturalistic narratives with annotated character states would test whether our synthetic results transfer to ecologically valid settings.
Tracking additional attribute types---beliefs, emotions, relationships---would determine whether the capacity limit is per-character or per-attribute-slot.
Finally, probing studies on open-weight models could directly measure the representational overlap that our behavioral results imply.
